\documentclass[UTF8,fontset=windows,zihao=-4,a4paper]{ctexart}
\usepackage[margin=22mm,top=25mm,bottom=23mm,headheight=16pt]{geometry}
\usepackage{graphicx,xcolor,amsmath,booktabs,tabularx,array,enumitem,fancyhdr}
\usepackage[unicode,colorlinks,linkcolor=teal!70!black,urlcolor=blue!45!black,pdfstartview=FitH]{hyperref}
\graphicspath{{figures/}}
\definecolor{ink}{HTML}{182C40}
\definecolor{accent}{HTML}{086C65}
\definecolor{muted}{HTML}{4B6171}
\definecolor{wash}{HTML}{EEF5F5}
\definecolor{warm}{HTML}{A04A23}
\hypersetup{pdftitle={从极端降水反演学科研},pdfsubject={YunTAPR-Net 案例讲义与科研方法}}
\ctexset{section={format=\Large\bfseries\color{ink},beforeskip=0pt,afterskip=12pt}}
\pagestyle{fancy}
\fancyhf{}
\fancyhead[L]{\small\color{muted}从极端降水反演学科研}
\fancyhead[R]{\small\color{muted}YunTAPR-Net}
\fancyfoot[L]{\small\hyperlink{nav}{目录}\quad\hyperlink{gl-data}{术语注释}\quad\hyperlink{sources}{证据与延伸阅读}}
\fancyfoot[R]{\small\color{muted}\thepage}
\renewcommand{\headrulewidth}{0.35pt}
\setlength{\parindent}{0pt}
\setlength{\parskip}{6pt}
\renewcommand{\arraystretch}{1.28}
\linespread{1.13}
\setlist{nosep,leftmargin=1.8em,topsep=4pt,itemsep=3pt}
\newcolumntype{L}{>{\raggedright\arraybackslash}X}
\newcommand{\term}[2]{\hyperlink{g-#1}{#2}}
\newcommand{\src}[1]{\hyperlink{s-#1}{\textsuperscript{[#1]}}}
\newcommand{\back}[2]{\hyperlink{#1}{\textcolor{muted}{回到#2}}}
\newcommand{\source}[3]{\hypertarget{s-#1}{}\textbf{[#1]}\enspace\href{#2}{#3}\par}
\newcommand{\chapterpage}[3]{\clearpage\hypertarget{#1}{}\section{#2}\textcolor{accent}{\bfseries #3}\par}
\newcommand{\sub}[1]{\par\medskip{\large\bfseries\color{ink}#1}\par}
\newcommand{\boxnote}[2]{\par\smallskip\noindent\colorbox{wash}{\begin{minipage}{\dimexpr\linewidth-2\fboxsep\relax}\vspace{3pt}\textbf{\textcolor{accent}{#1}}\par #2\vspace{3pt}\end{minipage}}\par\smallskip}
\newcommand{\diagram}[4][1]{\par\smallskip\noindent\makebox[\linewidth][c]{\includegraphics[width=#1\linewidth]{#2}}\par{\small\color{muted}\textbf{图 #3}\enspace #4\par}\smallskip}
\newcommand{\gloss}[4]{\par\hypertarget{g-#1}{}\textbf{\textcolor{ink}{#2}}\enspace #3\hfill{\small #4}\par}
\newcommand{\navline}[3]{\hyperlink{#1}{\textbf{#2}}\hfill\pageref{p-#1}\par{\small\color{muted}#3}\par\smallskip}
\newcommand{\mmh}{\mathrm{mm\,h^{-1}}}
\begin{document}
\thispagestyle{empty}
\vspace*{12mm}
{\large\bfseries\color{accent} 科研方法讲义\par}
\vspace{9mm}
{\fontsize{35}{45}\selectfont\bfseries\color{ink}从极端降水反演\\学科研\par}
\vspace{5mm}
{\Large\color{muted}以云南复杂地形 YunTAPR-Net 项目为例\par}
\vspace{8mm}
读完并完成练习后，检验自己能否独立写出一个可被数据驳倒的问题，设计公平比较，复算关键指标，并在组会上讲清楚证据能支持的结论。
\diagram{01-workflow.pdf}{1}{从问题、协议和数据到对照、诊断和表达。诊断发现的反例可以促使研究者重写问题。流程为教学示意。}
\boxnote{怎么使用这份讲义}{先读正文，点击青绿色术语可跳到定义和例子，注释有返回正文的链接。上标编号通往证据索引，便于核对原始公开记录。读完合上讲义，先做练习，再对照参考答案和组会评分表。}
\vfill
{\small\color{muted} 案例证据截至 2026 年 11 月 21 日 \hfill \hyperlink{nav}{进入目录 $\rightarrow$}\par}

\clearpage
\hypertarget{nav}{}
\section*{学习路径}
\navline{question}{1\quad 把研究兴趣写成可检验的问题}{交付问题卡，说明对象、比较和什么结果会改变判断。}
\navline{data}{2\quad 定义数据，再谈模型}{交付一张数据契约，区分观测时标、资料到达和参考窗口。}
\navline{model}{3\quad 把假说翻译为模型与对照}{用四条候选主张解释特征、网络和输出各在检验什么。}
\navline{experiment}{4\quad 设计比较并留下可复核的证据}{冻结配对样本、主要指标、保护线和停止规则。}
\navline{diagnosis}{5\quad 读懂强雨失效的三层证据}{复算检出和排序，理解条件分位数与固定点值。\enspace\hyperlink{freeze}{查看冻结协议实例}}
\navline{verification}{6\quad 从像元到暴雨过程}{写清同源参考、时间留出和独立观测各自的结论范围。}
\navline{communication}{7\quad 写结果、讲组会、回答追问}{输出一段结果、一页比较说明和三分钟口头讲解。}
\navline{exercises}{8\quad 练习与自评}{设计一项研究、复算数字、给结论降到证据允许的强度。}
\medskip
\hyperlink{gl-data}{数据与模型术语}\quad
\hyperlink{gl-eval}{统计与指标术语}\quad
\hyperlink{gl-method}{研究设计与输出术语}\par
\hyperlink{sources}{项目证据与原始阅读}
\vfill
\boxnote{案例范围}{讲义采用 Obsidian 项目记录与相应公开证据。项目当前科学状态为 \texttt{NOT\_VERIFIED}。文中明确区分实际结果、教学算例和拟议的评价步骤。\src{P1}}

\chapterpage{question}{把研究兴趣写成可检验的问题}{从“想做极端降水”到一张能够被反例推翻的问题卡}
\label{p-question}
云南山地的云顶辐射、低层水汽、近地风与地形相互作用，地面降水与卫星看到的云并非一一对应。YunTAPR-Net 研究的是复杂地形下的\term{retrieval}{卫星降水反演}，目标是用多源输入估计给定时间窗的降水；\term{latency}{近实时}适用性还依赖输入抵达时间与处理延迟。\src{P1}\src{P2}\src{R3}

\sub{一张可检验的问题卡}
\begin{tabularx}{\linewidth}{@{}p{3.2cm}L@{}}
\toprule
要素 & 用云南案例填写 \\
\midrule
研究对象 & 已配对的云南参考格点、半小时目标窗、高海拔强雨子集。\\
可操纵因素 & 固定网络及其余条件，比较真实动态地形输入与相同位置数值归零。\\
事前预期 & 真实地形输入会降低既定高海拔强雨子集的物理\term{pinball}{Pinball 损失}。\\
反驳条件 & 收益方向不稳定，或误报、弱雨等保护线失守。结论强度还取决于事件覆盖与重复性。\\
证据范围 & 先做\term{split}{开发集}比较，再检验不同\term{seed}{随机种子}、天气过程和时间留出。\\
\bottomrule
\end{tabularx}

\sub{读文献时，把结果追问到底}
以 Hydronn 原文为练习。先找研究区、卫星输入、训练参考与测试对象，再定位哪张表支持哪个结论。最后记录要迁移的设计，以及云南的地形、时间支撑或参考来源可能改变什么。Obsidian 文献索引可沿“静止卫星反演→多源融合→概率与强雨→实时订正”这条路线继续读。\src{R3}

\sub{先写下最容易混淆的解释}
如果高海拔只在两场相似天气中出现较多，局部收益可能反映过程构成。如果归零改变了网络参数或训练样本，差异也可能来自容量或数据。读完一篇文献后，写一条“可迁移的做法”和一条“需要重新验证的主张”，并标出区分这两种解释的比较。\src{P4}\src{P6}

\boxnote{把方法迁移到自己的课题}{用一句话填空。「在［人群或区域、时间和单位］上，把［变量］改为［候选］，相对［对照］，预计［主指标］改变，且［保护指标］达标。若［观测结果］出现，我会收窄或放弃主张。」}

\chapterpage{data}{定义数据，再谈模型}{一条样本要说得清哪里、何时、用的是什么版本}
\label{p-data}
\begin{tabularx}{\linewidth}{@{}p{2.65cm}L@{}}
\toprule
数据身份 & 本项目案例中的定义 \\
\midrule
目标与单位 & \term{imerg}{IMERG V07 Final} 半小时平均雨强。设起点 \(T\)，目标窗 \([T,T+30\,\mathrm{min})\)，分析时刻 \(A=T+30\,\mathrm{min}\)，时间记为 \term{time}{UTC}。 \\
卫星输入 & 一个记录配置为七个\term{satellite}{光谱通道}、六个名义卫星时槽 \(A-60,A-50,A-40,A-30,A-20,A-10\) 分钟，解码时观察结束时刻须满足协议上的因果限制。\\
环境与地形 & \term{gfs}{GFS f006} 提供回顾性环境输入。地形来自 \term{dem}{DEM}。环境字段的变量身份、单位、气压层和\term{mask}{有效掩膜}写进输入契约。 \\
空间与划分 & \term{grid}{SP04} 约 \(0.1^\circ\) 网格及云南监督掩膜。案例以 2023 训练、2024 开发，2025 最终测试期封存。\\
\bottomrule
\end{tabularx}
\src{P1}\src{P2}\src{P3}

\diagram{02-timeline.pdf}{2}{通用时间示意。图中时间关系仅用于解释资料何时可用。项目的目标窗在分析时刻 \(A\) 结束，部分输入名义时槽可位于目标窗内。}

\textbf{三项开跑前检查。} 第一，逐条核对输入观测起止、模式起报、预报有效和档案实际可用时间；回顾性资料对齐为核查前几项绑定提供依据。第二，检查经纬度轴方向、格点中心、范围、单位、掩膜与陆地以下气压层。第三，先划分再拟合\term{scaler}{标准化}统计，构造同窗同格点的两臂共同有效样本。\src{P1}\src{P2}\src{P3}

\textbf{核查实例。} 若卫星文件标记 \(A-10\) 分钟，实际扫描结束晚于 \(A\)，就不能只看文件名入选。若某气压层在山区地表以下，按项目契约标记无效；若混入 Pa 和 hPa，错误应在标准化前被发现。\src{P2}\src{P3}

\boxnote{半小时换算}{若完整半小时累计雨量是 \(H\) mm，则对应平均雨强为 \(2H\) mm/h。\(H=10\) mm 恰好等于 20 mm/h。项目把强雨定义为严格 \(>20\) mm/h，因此恰好等于阈值的值归入非强雨。}

\chapterpage{model}{把假说翻译为模型与对照}{四个方向分别提出四个可以被检验的问题}
\label{p-model}
\begin{tabularx}{\linewidth}{@{}p{3.15cm}LL@{}}
\toprule
项目方向 & 为什么值得试 & 怎样识别贡献 \\
\midrule
多元信息协同 & \term{satellite}{多时相、多光谱}补充云演变及通道差异，环境和地形给出位置与状态背景。 & 用\term{ablation}{消融}分别固定通道或时间槽，其余比较协议与样本身份一致。 \\
动态地形强迫 & \term{dote}{风场与地形梯度}可提供上坡运动的近似表征。 & REAL/ZERO 把对应特征在同架构中归零，配对核对真值、掩膜、海拔和雨强分层。独立暴雨事件数仍未知。 \\
环境物理约束 & 当前\term{mee}{MEE 环境先验}有 13 个变量，并绑定物理量、单位与地下层掩码。 & 在同人口配对中检验环境增量。当前为条件输入，硬守恒及业务可用性分别需要证据。 \\
U-Net 双头 & \term{unet}{U-Net}融合空间上下文。\term{heads}{双头}分别输出有雨概率与 32 个单调条件分位数。 & 分开评估发生、条件强度与强雨越线。固定\term{readout}{读出}后再比较阈值检出。 \\
\bottomrule
\end{tabularx}
\src{P2}\src{P3}\src{P4}\src{P6}

\sub{从数学式读出不同的预测}
令 \(R\) 是雨强，\(W=\{R>0.1\,\mmh\}\)，\(x\) 是输入。发生头估计 \(p(x)=P(W\mid x)\)，条件头估计 \(Q_\tau(R\mid W,x)\)。两者可以共同描述强雨风险。对 \(u=20\,\mmh\)，有
\[
P(R>u\mid x)=p(x)\,P(R>u\mid W,x).
\]
这个等式依赖 \(\{R>u\}\subset W\)。项目用有限节点表达\term{quantile}{条件分位数}，节点之间与两端的分布还需额外假设才能给出单一越线概率。q32 的 \(\tau=0.984375\) 表示条件分位水平。\src{P2}\src{P7}

\sub{让机制主张承担可被推翻的预期}
若上坡气流表征是有效信息，合理预期是同样雨强下，收益随地形坡向与风向的相互关系改变。对应检验需要事前定义风向、坡向和海拔分组，并把 REAL/ZERO 差值在多个天气过程中重复比较。当前已有海拔分层差异，风向机制仍需专门的证据。\src{P3}\src{P6}

\boxnote{最小算例}{教学例。若两个模型的有雨概率同为 0.8，而条件下超过 20 的概率分别为 0.05 与 0.25，则强雨概率分别为 0.04 与 0.20。仅看发生头会错过这个差异。}

\chapterpage{experiment}{设计比较并留下可复核的证据}{先写比较规则，再看结果}
\label{p-experiment}
\begin{enumerate}
\item \textbf{冻结因变量。} 写明指标、单位、分母、样本身份、比较方向与\term{guardrail}{保护线}。例如事前同时要求强雨损失下降、弱雨和误报保持在容许范围。
\item \textbf{冻结对照。} 绑定训练和评价的\term{population}{样本集合}、数据版本、顺序、初始化、学习率调度、网络结构、输出规则与\term{checkpoint}{固定检查点}。记录唯一有意改变的因素。
\item \textbf{先用小检查排错。} 用人工张量试验临界值、缺帧、无雨、错误单位和网格错位。先查输入契约、算分器和账本，再用获准的小真实样本检查整条路径。
\item \textbf{保存可复核材料。} 记录代码提交、配置、随机种子、检查点指纹、样本/掩膜身份、有效计数、停止原因和失败运行。一次恢复的重放和一个独立 seed 各自记账。
\item \textbf{报告差值与不确定性。} 同一场景两臂相减。复现时分开报告种子差异和暴雨过程差异。若采用\term{block}{事件级重采样}，先说明区块构造和有效独立性。
\end{enumerate}

\diagram{03-evidence.pdf}{3}{工程核验、开发比较、时间留出与独立观测对应不同的证据问题。箭头示意证据范围扩展，图不代表项目已通过这些阶段。}

\textbf{历史案例。} B3/B4 三个 seed 在同一 2024 Dev 上的 Brier、AUROC、AP 和条件 Pinball 总体同向改善，q32 方向存在不一致，强雨仍有退化。B0/B1 的历史训练和评价人口与 B3/B4 不同，跨人口拼排无法识别多光谱增益。\src{P4}

\textbf{为何要保留代价？} COSINE9 把学习率时程从原方案缩至九轮。与固定对照相比，强雨 log Pinball 约降 10.00\%，弱雨 log Pinball 却升约 7.59\%。二者在不同分层、不同分母上统计。先写清样本构成与保护线，再判断变更是否值得保留。\src{P5}\src{P12}

\boxnote{复现前先自问}{别人拿到这份记录，能否找到同一代码和检查点、重建同一像元的配对、复算同一个分母、识别什么时候这个候选已经失败？}

\chapterpage{diagnosis}{读懂强雨失效的三层证据}{先问指标在评估哪个输出，再决定改模型还是改读出}
\label{p-diagnosis}
\diagram[0.83]{04-diagnosis.pdf}{4}{发生、条件幅值、尾部识别要分别回答。第三条轨道还应区分节点越线、概率排序与固定点值检出。}

\textbf{发生层。} \term{brier}{Brier} 评估有雨概率与真实有雨指示的平方差。\term{rank}{AP、AUROC}关注分数排序。若正例定义为有雨，指标只回答有雨问题；AP20、AUROC20 则把强雨作为正例。\src{P2}\src{P8}\src{R6}

\textbf{条件幅值层。} \term{pinball}{Pinball} 衡量指定分位点的误差。REAL 相对 ZERO 在同一 seed、固定 E9、2024 Dev 的高海拔强雨子集上，物理 Pinball/20 降 0.1367\%，log Pinball 降 0.2106\%。该子集含 178 个\term{event}{像元时}，落在 80 个场景内。两个较低海拔层方向相反，高海拔 \(>50\,\mmh\) 样本为零。全域 log 损失略升，Brier 略降，AP 和 AUROC 略降。两臂三种点估计均命中 0/717，高海拔 q32 覆盖由 6/178 增至 11/178。该比较支持单 seed 的局部描述。\src{P6}

\textbf{尾部识别层。} \term{tail}{q32 覆盖与节点越线}各有自己的判断条件。TAIL20 原九项筛查通过八项。强雨物理 Pinball/20 从 0.506865 降到 0.504120，真强雨 q32 越线 160 增至 179。唯一失守项是非强雨 q32 越线，5,769 增至 6,602。A 点值仍命中 0/717，候选封存。\src{P7}

\sub{A20 的命中增加，要怎样解释？}
同一批 963 个场景、3,303,090 个有效像元时中，强雨正例 717 个。下表评价同一固定点值 A 的阈值检出和跨阈值排序。\src{P8}\par
\begin{tabularx}{\linewidth}{@{}Lrr@{}}
\toprule
指标 & COSINE9 & A20 \\
\midrule
固定 \(>20\,\mmh\) 的\term{confusion}{命中 / 误报} & 0 / 0 & 627 / 177,098 \\
全体 AP20 & 0.018691 & 0.007786 \\
全体 AUROC20 & 0.967601 & 0.968501 \\
海拔 \(\ge 2{,}000\) m AP20 & 0.008314 & 0.004929 \\
海拔 \(\ge 2{,}000\) m AUROC20 & 0.978222 & 0.971738 \\
\bottomrule
\end{tabularx}

A20 命中率约 87.45\%，报出的强雨中约 0.35\% 是真强雨。全体 AUROC20 略升，但两个已报告范围的 AP20 均下降。原筛查九项仅过两项，固定阈值误报线失守。排序补证后按预设规则结束该候选。点值排序在这里仍需和\term{calibration}{概率校准}分开识别。\src{P8}\src{P13}\par\hyperlink{case-figures}{打开两幅结果图 $\rightarrow$}

\clearpage
\hypertarget{case-figures}{}
\section*{用两幅结果图练习判读}
{\color{accent}\bfseries 先找比较器、样本量和基准线，再读差值方向。}\par

\diagram{05-terrain.pdf}{5}{REAL 相对 ZERO 的强雨分位损失变化。负值表示损失下降。三个海拔层的 \(n\) 都是像元时数。固定 B7/seed2029/E9 的 2024 开发配对，未给出独立天气过程数或跨种子区间。\src{P6}}
\textbf{读图动作。} 先按同一指标横向看三个海拔层，再比较同一层的两个损失尺度。选定高海拔子集后，句子里保留中低海拔的反例与高海拔极端样本不足的限制。\par

\diagram{06-ranking.pdf}{6}{同一强雨定义下的 A 点值 AP20。参考为 COSINE9，候选为 A20。全体包括高海拔子集，不能把两行视作独立重复。阈值的命中和误报见\hyperlink{diagnosis}{第5节}。\src{P8}}
\textbf{读图动作。} 问听众为何 A20 在固定阈值命中更多的同时，AP20 仍下降。如果只给出 AP20 图，请补分母、固定阈值计数及原先接受的误报线。\par
\textbf{补齐原筛查。} A20 两项通过的是 A 的命中与真强雨 q32 越线。强雨物理 Pinball/20 从 0.506865 升至 0.628274，非强雨 q32 越线从 5,769 增至 352,680，弱雨损失和发生指标也退化。\src{P13}\par



\clearpage
\hypertarget{freeze}{}
\section*{把门槛写成可执行的协议}
\label{p-freeze}
{\color{accent}\bfseries 2026-11-21 的 IMERG 开发评价冻结实例}\par

\textbf{比较器与身份。} B7/REAL/COSINE9，seed2029、固定 E9。共同有效监督像元时 3,303,090，参考有雨 407,820，严格大于 20 的强雨 717，弱雨 \(0.1<R\le1\,\mmh\) 共有 246,473。先核对真值、掩膜、海拔分区和概率节点合同，再算差值。\src{P10}

\begin{tabularx}{\linewidth}{@{}L L@{}}
\toprule
事前要求 & 它约束的失败模式 \\
\midrule
强雨物理 \term{pinball}{Pinball/20} 严格下降 & 防止只把少量输出推过阈值，却没有改善条件幅值。\\
真强雨 q32 越线 \(\ge160\)，非强雨 q32 越线 \(\le5{,}769\) & 防止单纯整体提高上尾，同时放大非强雨位置的尾部支撑。\\
固定 A 的命中 \(>0\)，误报 \(=0\) & 要求固定读出恢复检出，同时守住当前冻结的误报线。\\
弱雨 log Pinball、发生 Brier 不升，发生 AP、AUROC 不降 & 检查候选没有通过牺牲较常见降水与发生表现取得局部改善。\\
\bottomrule
\end{tabularx}

\textbf{实际算分顺序。} 对强雨真值样本，物理 Pinball/20 先在 32 个物理分位上平均，再在强雨像元时上平均。固定 \term{readout}{A 点值}是发生概率乘以 32 个物理分位数的算术平均，以 64 位浮点计算，转为 32 位浮点后按严格 \(A>20\) 判定。AP20 用同一 A 分数，是全域和高海拔都需报告的排序信息，在本次新协议中属于辅助报告。历史 A20 按该实验的预设排序停止规则封存。\src{P8}\src{P10}

\boxnote{失败也要留下什么}{记录尝试改了什么、为何可能有效、哪项结果推翻预期、哪些结果仍成立，以及下一项需要什么证据。当前冻结了开发筛查和最终测试边界，待核实连续事件覆盖；截至本快照尚无准入的新候选，2025 的实际测试范围与执行仍待研究者绑定。\src{P10}}

\sub{把实验阶段看作逐项排除解释}
B0/B1 比较时序，B2/B3 比较多光谱输入条件，B3/B4 检查 GFS 环境增量，随后 B5/B6 引入静态地形，B7/B8 探索动态地形与融合方式。各段有自己的比较器、人口和种子。列出这条路线时，逐项报告实际匹配的比较和各自反例。\src{P4}\src{P11}


\chapterpage{verification}{从像元到暴雨过程}{先确认参考口径，再扩大验证范围}
\label{p-verification}
\sub{同源参考和时间留出分别回答什么？}
当前采用 \term{reference}{IMERG V07 Final 同源参考}。它也提供学习标签，因此配对结果衡量对该参考的表现。最终候选和协议冻结后，2025 可用于同源时间留出检验；截至案例日期仍封存。IMERG 的 V07 是版本，Final Run 是处理流，2025 Final Test 是数据划分。\src{P1}\src{R1}

\sub{为什么 717 个像元时没有给出 717 场暴雨？}
像元时是一个目标窗内的一个像元；场景是一次格点场采样；天气过程跨多个相邻时刻和空间位置。2024 Dev 的 717 个强雨像元时分布于 207 个含强雨窗口，多数相邻场景约隔 6 小时。仅凭稀疏窗口无法恢复完整暴雨过程，也不能把相邻像元视作独立样本。\src{P1}

事件评价要先用连续参考建立\term{event}{过程目录}，事前写清邻域、时间槽、位移容忍、最小面积、持续时间、缺测和截断。再映射共同有效的模型窗，分别报告目录过程数、可匹配数、可评价数及匹配覆盖。当前协议为草案，覆盖待核实。若仅能核对稀疏时窗，就把结论限定在匹配窗内。\src{P1}

\sub{什么能增强独立物理验证？}
站点给出点测，雷达 \term{qpe}{QPE} 是经过估计的面积结果，\term{dpr}{DPR} 是卫星过境时的瞬时足迹。纳入前检查来源链与标签是否共享观测、时间支撑、空间代表性、地形遮挡、版本和缺测。现行路线已确认 IMERG 参考并冻结开发评价，站点、雷达和 DPR 为可选补充，独立观测证据仍为 \texttt{PARTIAL}。\src{P1}\src{P10}\src{R1}

\sub{怎样报告不确定性}
训练随机性用同规则的配对 seed 差值描述。天气差异用可核对的过程或时块来表述。若进行重采样，两臂在同一抽样中的事件集合要一致。报告覆盖、样本数、差值区间和正负方向。条件不满足时保留描述性效应，说明哪项不确定性仍未估计。

\boxnote{检验一条结论的适用范围}{问自己「换年份、换天气过程、换观测来源」三件事是否都实际做过。每多写一个适用范围，都要指出相应比较的样本、协议与结果。}

\chapterpage{communication}{写结果、讲组会、回答追问}{把研究结论写成别人能逐项核查的句子}
\label{p-communication}
\sub{可直接改写成结果段的案例}
「在 2024 Dev 的 963 个配对场景、3,303,090 个共同有效像元时中，比较固定 E9 的 COSINE9 与 A20 点值。A20 在严格 \(>20\,\mmh\) 的阈值下命中 627/717 个强雨像元时，并产生 177,098 次误报。全体 AP20 从 0.018691 降至 0.007786，高海拔子集同向下降，因此按事前准则封存该候选。这一比较衡量所用开发场景对 IMERG 参考的点值检出和排序。」\src{P1}\src{P8}

\sub{三分钟组会分配}
\begin{tabularx}{\linewidth}{@{}p{2.15cm}L@{}}
\toprule
用时 & 要讲出的句子 \\
\midrule
0:00--0:35 & 我们想估计云南山区半小时降水。强雨点值长期难以过阈值，因此这一轮要分辨幅值不足和越线后误报。\\
0:35--1:15 & 两臂使用相同场景、真值与掩膜，我用固定该阈值的命中和误报检验可用性，再用 AP20 看强雨在分数中的排序。\\
1:15--2:20 & 报告参考臂、候选、分母和误报。指出 AP20 下降与 AUROC20 小幅上升可以同时发生，并说明固定阈值验收为何失败。\\
2:20--3:00 & 这仍是对 2024 Dev 上同源参考的比较。下一份能够改变结论的材料是事前定义的、覆盖可核验的过程评价或留出结果。\\
\bottomrule
\end{tabularx}

\sub{给听众留一份复核记录}
在图注或一页说明上，保留\term{claim}{主张与证据}、输入和参考版本、时空窗、两臂唯一差异、阈值、分母、效应方向、反例与局限。附可定位的公开收据。写“改善”时附绝对数和变化方向，写“显著”前给出推断方法、独立性假设及区间。

\sub{一页说明写什么}
\textbf{标题}用一个问题。左侧写对照身份、实际改变和评价范围；右侧放带分母的结果，至少保留一项反例。末行写当前结论、可推翻它的下一项观测，以及供听众查证的收据编号。\back{sources}{证据索引}

\boxnote{组会常见追问}{「为什么不用 2025 选个更好的阈值？」开发选择必须和最后一次检验隔离。「为什么有雨概率高还漏强雨？」发生与条件强度是不同输出，固定 A 还取决于条件均值。「为什么现有结果不能证明地形机制？」因为目前局部变化、反例与广泛复现的证据范围不同。}

\chapterpage{exercises}{练习与自评}{闭卷完成，再翻到参考答案}
\label{p-exercises}
\sub{练习一　设计一项可被否证的研究}
选一类环境字段或一个新增卫星通道，写 150 字问题卡。交付比较臂、样本身份、主指标、至少两条保护线、失败条件，以及失败之后需要核对的一个原因。然后把同一模板迁移到你的课题。

\sub{练习二　审查时间和单位}
教学例：目标窗为 12:00--12:30 UTC。一场记录累计 11 mm；某个有效时刻为 12:30 的模式产品在 12:42 才可下载，而业务拟在 12:35 发布。写出平均雨强、是否为强雨，以及该资料能否用于这次业务发布。说明时窗完整性有何影响。

\sub{练习三　算出研究结论的代价}
用案例中的 717 个强雨像元时、A20 的 627 命中和 177,098 误报，求漏报数、POD 和精确率。对照全体 AP20，写一句 60 字以内的结果，带上数据范围和停止理由。

\sub{练习四　把一段过强的结果写回证据}
改写「加入动态地形后，模型能准确检测山区所有暴雨，已可用于预警。」必须引用海拔层差异、像元时与事件的区别、固定点值结果，并留下业务可用性条件。随后找同学当评委，完成三分钟讲解。

\bigskip
\hyperlink{answers}{核对参考答案 $\rightarrow$}

\clearpage
\hypertarget{answers}{}
\section*{参考答案与输出评分}
\textbf{一。} 示例：在固定共同开发样本和固定检查点比较「原环境输入」与「只新增一个水汽字段」。以参考强雨物理 Pinball 为主指标，同时保护弱雨分位损失与固定阈值误报。若主指标无改善或任一保护线越界，结束该候选，优先查单位和地下层掩膜。具体容差需在看结果前确定。

\textbf{二。} 11 mm 除以 0.5 h 等于 22 mm/h，满足严格 \(>20\)。产品 12:42 才到达，晚于 12:35，不符合这次发布。若累计窗覆盖不完整，标记缺测并核对资格，不能按半小时直接换算。

\textbf{三。} 漏报 \(717-627=90\)，POD \(=627/717\approx87.45\%\)，精确率 \(=627/(627+177098)\approx0.35\%\)。示例：「2024 开发配对中，A20 虽命中 627/717，误报 177,098 且 AP20 下降，未通过冻结准则。」对照臂无正预测时，精确率的分母为零，报告未定义或注明工具的约定。

\textbf{四。} 示例：「2024 Dev 的单 seed 配对中，REAL 在高海拔强雨像元时的物理分位损失略降，两个较低海拔层退化。两臂固定点估计在强雨阈值下均零命中，事件覆盖和业务可用性仍待核实。」\src{P6}

\sub{同伴评分表}
\begin{tabularx}{\linewidth}{@{}p{3.4cm}L@{}}
\toprule
维度 & 满分 2 分的可观察表现 \\
\midrule
问题与设计 & 能给出可被推翻的主张、匹配对照和事前保护线。 \\
数据与时空 & 能说清参考、目标窗、单位、共同样本和资料可用时间。 \\
读数与解释 & 能复算分母，区分分位损失、点值检出、概率和排序。\\
表达与追溯 & 能讲清结论边界和反例，并指出两份可核对证据。\\
\bottomrule
\end{tabularx}
每项可独立完成记 2 分，提示后完成记 1 分，无法完成记 0 分。教学建议是 7/8 分且数据时空、读数两项均满分后，再尝试独立讲解。选一条自己的研究记录，交出一页问题卡、一段结果和一次三分钟讲解作为迁移测验。

\clearpage
\hypertarget{gl-data}{}
\section*{链接注释 · 数据与模型}
{\small 词条后的数字通往来源，返回链接指向对应正文。另读 \hyperlink{gl-eval}{统计与指标}、\hyperlink{gl-method}{研究设计与输出}。}

\gloss{retrieval}{卫星降水反演}{用卫星测到的辐射等信息推断对应时空范围的降水。未来时段预报还要定义预报提前量与未来目标。}{\src{P2}\src{R3}\back{question}{第1节}}
\gloss{latency}{近实时与时延}{从观测结束到可用产品输出的耗时应可核验，涉及数据到达、处理、缺报和发布。历史有效时间匹配不能单独证明业务时效。}{\src{P2}\back{question}{第1节}}
\gloss{imerg}{IMERG V07 Final}{NASA 的 GPM 多卫星融合降水产品。V07 是版本，Final 是事后处理流，带月尺度雨量站资料调整。参考产品存在自身误差。}{\src{R1}\src{P1}\back{data}{第2节}}
\gloss{time}{时间支撑与 UTC}{时间支撑表示测量实际覆盖的区间。UTC 是统一时间基准。目标窗左右端是否纳入、观测有效时间与文件获取时间都应记录。}{\src{P2}\back{data}{第2节}}
\gloss{satellite}{光谱通道、亮温与时槽}{通道对应不同光谱测量，亮温是将辐射换算成等效黑体温度。时槽是名义采样位置，真正的观测起止仍需解码核对。云顶信号与地面雨强的关系依赖天气环境。}{\src{R3}\src{P2}\back{data}{第2节}}
\gloss{gfs}{GFS、f006 与资料版本时效}{GFS 是全球数值天气模式。f006 表示距起报 6 小时，预报有效时刻通常为起报加时效。可用时间还要核查生产与传输，因此需要实际可用档案版本。}{\src{R2}\src{P3}\back{data}{第2节}}
\gloss{grid}{格点与 SP04}{SP04 为项目的空间绑定协议。复算需匹配坐标参考、经纬轴、像元位置及监督范围。经纬度网格的同度数面积随纬度变。}{\src{P1}\src{P2}\back{data}{第2节}}
\gloss{mask}{掩膜}{布尔资格标记。云南监督掩膜定空间范围，数据质量掩膜排除无效值，地下层掩膜排除对应地表以下的气压层。比较使用共同有效格点。}{\src{P3}\back{data}{第2节}}
\gloss{scaler}{标准化}{常见写法 \(x'=(x-\mu)/\sigma\)。统计量用训练数据拟合，按同一字段、单位和身份应用到后续划分，避免开发或测试信息进入训练。}{\src{P2}\src{P3}\back{data}{第2节}}
\gloss{dem}{DEM 与地形梯度}{DEM 是海拔栅格。以水平米制距离求 \(\nabla h=(\partial h/\partial x,\partial h/\partial y)\)，单位为 m/m。应记录投影、聚合、缺失与接缝。}{\src{P3}\back{data}{第2节}}
\gloss{dote}{动态地形强迫与 DOTE}{可用 \(u\,\partial h/\partial x+v\,\partial h/\partial y\) 理解沿地形的运动表征，风用 m/s 时结果单位为 m/s。正值表示沿所定义坐标吹向上坡。项目特征还绑定物理场与掩膜，具体实现见工程记录。}{\src{P3}\src{P6}\back{model}{第3节}}
\gloss{mee}{MEE 与物理条件化}{项目环境分支把 GFS 物理字段作为条件输入，并验证变量、气压层、单位和掩膜。显式守恒约束应说明守恒方程、离散方式与检验。}{\src{P3}\src{P4}\back{model}{第3节}}
\gloss{unet}{U-Net 与跳接}{编码阶段聚合上下文，解码阶段恢复空间表达。跳接把相应尺度的空间信息传至解码层。项目借用其结构思想做格点估计。}{\src{R4}\back{model}{第3节}}
\gloss{heads}{发生头与条件分位头}{共用特征后分别输出有雨概率与有雨条件下的分位节点。条件头在参考有雨样本上直接接受分位监督。}{\src{P2}\back{model}{第3节}}

\clearpage
\hypertarget{gl-eval}{}
\section*{链接注释 · 统计与指标}
\gloss{quantile}{分位数与条件分位数}{\(Q_\tau=\inf\{r:F(r)\ge\tau\}\)。条件分位数使用指定条件下的分布。本项目 \(\tau_i=(i-0.5)/32\)，最高为 0.984375。例：若模型校准良好，按同样条件汇总时，约 90\% 的观测应不超过预测的 0.9 分位数，离散值的并列需要注明。}{\src{P2}\src{R5}\back{model}{第3节}}
\gloss{pinball}{Pinball 损失}{令 \(e=y-q_\tau\)，则 \(\rho_\tau(e)=e(\tau-\mathbf1[e<0])\)。在 \(\tau=0.9\) 时，低估 2 单位得 1.8，高估 2 单位得 0.2。案例同时报告 \(\log(1+R)\) 域和物理域。物理 Pinball/20 按引用报告的固定尺度归一化。报告样本、分位水平、尺度与分母。}{\src{P2}\src{P7}\src{R5}\back{diagnosis}{第5节}}
\gloss{brier}{Brier 分数}{对样本计算 \((p_i-z_i)^2\) 后取平均，\(z_i=1\) 表示事件发生。分数越低越好。例如预测 0.8，若事件发生得 0.04；若没发生得 0.64。}{\src{R6}\back{diagnosis}{第5节}}
\gloss{rank}{AP、AUROC 与正例定义}{AP 按召回增加加权汇总精确率。AUROC 汇总阈值变化时真阳性率与假阳性率的关系，常可解释为正例得分超过负例的概率，并列按半计。强雨稀少时，报告正例比例和 AP，并结合固定阈值计数。换正例定义或样本会改变可比性。}{\src{P8}\src{R7}\back{diagnosis}{第5节}}
\gloss{confusion}{命中、误报、漏报和阈值指标}{TP 为真强雨且报强雨，FP 为非强雨却报强雨，FN 为真强雨却没报。\(POD=TP/(TP+FN)\)，精确率 \(=TP/(TP+FP)\)，误报比 \(=FP/(TP+FP)\)，假阳性率 \(=FP/(FP+TN)\)。分母为零时明确报告处理约定。}{\src{P8}\back{diagnosis}{第5节}}
\gloss{tail}{覆盖、节点越线与尾部识别}{覆盖比较真值 \(R\) 与节点 \(q_\tau\)，例如计 \(R\le q_{32}\)；节点越线比较 \(q_{32}>20\)。固定点值检出则比较读出 A 与 20。有限节点的兼容概率界限反映信息不足，抽样置信区间反映估计不确定性。}{\src{P6}\src{P7}\back{diagnosis}{第5节}}
\gloss{readout}{读出、A 点值与 \(P20\)}{读出把模型输出转成可比较的量。固定 A 为发生概率乘以32 个物理分位节点的算术平均。若有雨定义 \(R>0.1\)，精确的全分布均值还会含低雨强条件项。\(P20\) 指超过 20 的概率，与 A 的物理雨强值有不同单位和含义。}{\src{P2}\src{P7}\back{model}{第3节}}
\gloss{calibration}{概率校准}{在概率近似 0.2 的一组独立可比样本中，实际发生率应接近 0.2。排序相同也可有不同校准。概率可靠性需要足量样本、置信信息及清晰的拟合与评价划分。}{\src{R6}\src{P7}\back{diagnosis}{第5节}}
\gloss{block}{配对与区块重采样}{先在同一单位上做候选减对照，再按事前确定的事件或时间块取样以保留块内相关性。区块过少、块间强依赖或覆盖偏倚时，明确推断受限。}{\src{P1}\back{experiment}{第4节}}
\gloss{event}{像元时、场景与天气过程}{像元时是格点乘目标窗，场景是一整个匹配格点场。过程是满足事前时空连通规则的一组降水暴露。应区分完整、匹配和可评价的过程。}{\src{P1}\src{P6}\back{verification}{第6节}}

\clearpage
\hypertarget{gl-method}{}
\section*{链接注释 · 研究设计与输出}
\gloss{split}{Train、Dev、Final Test 与泄漏}{Train 拟合参数，Dev 支持选择与诊断，Final Test 用于冻结方案后的评估。使用最终测试指标来挑候选、阈值或数据处理规则，会把测试信息泄露给开发。时间留出检验仍受参考来源限制。}{\src{P1}\src{P2}\back{question}{第1节}}
\gloss{seed}{随机种子}{控制初始化、抽样和顺序等随机过程。相同 seed 有助于配对，多个 seed 描述训练随机性。恢复一次运行不增加独立种子，多个随机种子也不增加独立天气过程。}{\src{P4}\back{question}{第1节}}
\gloss{ablation}{消融与匹配对照}{有意去掉或替换某种输入、模块或损失，用相同评价协议比较。若同步改变容量、数据或训练时长，应明确结果混入了这些因素。真实/归零对照是输入消融的一种。}{\src{P4}\src{P6}\back{model}{第3节}}
\gloss{population}{样本集合与共同有效样本}{比较所用样本必须有可核验的时空身份。相同场景数并不自动表示真值、格点、时标和掩膜相同。先核对集合，再比较指标。}{\src{P1}\src{P8}\back{experiment}{第4节}}
\gloss{guardrail}{主指标、保护线与停止规则}{主指标回应研究问题。保护线限定可接受代价，如弱雨误差或误报。停止规则规定候选何时封存或回到诊断，阈值应在见到结果前写定。}{\src{P7}\src{P8}\back{experiment}{第4节}}
\gloss{checkpoint}{检查点、E9 与指纹}{检查点保存某一训练状态。E9 是第九轮的固定终点，和事后挑最佳轮次的规则不同。文件指纹用于确认身份，还需数据与评估记录才能回溯结果。}{\src{P4}\src{P5}\back{experiment}{第4节}}
\gloss{reference}{参考一致性与独立验证}{同源参考评价沿用与训练标签相同的产品来源。独立验证要求检查待验证模型与观测的来源链、共享订正与覆盖，并在可比时空支撑上配对。}{\src{P1}\src{R1}\back{verification}{第6节}}
\gloss{qpe}{雷达 QPE 与站点代表性}{QPE 是定量降水估计。雷达反射率到地面雨强可能受遮挡和订正影响，站点的点测与格点面积平均存在尺度差异。}{\src{P9}\back{verification}{第6节}}
\gloss{dpr}{GPM DPR}{双频降水雷达。过境足迹能提供局部观测，瞬时空间足迹与半小时面积平均须按评价协议处理，并核对资料源链。}{\src{P9}\src{R1}\back{verification}{第6节}}
\gloss{claim}{主张、结果与科学验收}{结果是某口径下的观测差值，主张增加了范围和解释。工程完成、开发筛查、机制归因与泛化各有证据条件。正文应同时给出关键反例。}{\src{P1}\src{P4}\back{communication}{第7节}}

\clearpage
\hypertarget{sources}{}
\section*{证据与延伸阅读}
\newcommand{\repo}{https://github.com/flyeagle-cex/YunTAPR-Net/blob/23900b20579afb408727f9a7f2aaf0b4d286f712}
{\small 讲义从 Obsidian 的项目主页、研究设计与证据地图、四项创新、山区强雨、TAIL20、A20、IMERG 参考评价快照及文献索引组织而成。下列项目链接指向可复核的固定提交和公开资料。}

\sub{项目记录}
\source{P1}{\repo/docs/research_readiness/20261121_stage45_observation_feasibility/IMERG_REFERENCE_EVENT_PROTOCOL_V1.md}{IMERG 参考评价与事件协议 V1}
\source{P2}{\repo/docs/phase_b_v2_protocol_candidates/v1_20261010/METHODS_AND_EXPERIMENTAL_DESIGN.md}{数据时间绑定、训练划分和概率输出协议}
\source{P3}{\repo/docs/research_readiness/20261115_core_innovation_repair/OBSIDIAN_SNAPSHOT.md}{四项创新主线与物理输入修复快照}
\source{P4}{\repo/docs/research_readiness/20261025/B0_B4_STAGE_ACCEPTANCE.md}{B0--B4 阶段交接与匹配比较}
\source{P5}{\repo/docs/research_readiness/20261110_cosine9/PROTOCOL.md}{COSINE9 固定九轮比较协议}
\source{P6}{\repo/docs/research_readiness/20261115_mountain_pair_results/OBSIDIAN_TERMINAL_SNAPSHOT.md}{REAL/ZERO 山区强雨配对结果}
\source{P7}{\repo/docs/research_readiness/20261119_conditional_intensity_diagnostic/STRONG_RAIN_DETECTION_EVIDENCE.md}{TAIL20 与强雨检出瓶颈}
\source{P8}{\repo/docs/research_readiness/20261119_conditional_intensity_diagnostic/A20_POINT_RANK_RESULTS.md}{A20 与 COSINE9 配对排序结果}
\source{P9}{\repo/docs/research_readiness/20261121_stage45_observation_feasibility/STAGE45_EVENT_MATCHING.md}{独立观测与事件配对的可行性口径}

\source{P10}{\repo/docs/research_readiness/20261121_stage45_observation_feasibility/STAGE2_EVALUATION_FREEZE_V1.md}{冻结的开发评价、指标与最终测试边界}
\source{P11}{\repo/STAGE_ACCEPTANCE.md}{地形分支与各试验阶段验收}

\source{P12}{\repo/docs/research_readiness/20261110_cosine9/SAVED_TRADEOFF.json}{COSINE9 的雨强分层代价收据}

\source{P13}{\repo/docs/research_readiness/20261119_conditional_intensity_diagnostic/A20_RESULTS.md}{A20 原固定阈值验收与九项筛查}

\sub{专业知识原始来源}
\source{R1}{https://gpm.nasa.gov/sites/default/files/2023-07/IMERG_TechnicalDocumentation_final_230713.pdf}{NASA · IMERG V07 技术文档}
\source{R2}{https://emc.ncep.noaa.gov/emc/pages/numerical_forecast_systems/gfs.php}{NOAA/NCEP · Global Forecast System}
\source{R3}{https://amt.copernicus.org/articles/15/6907/2022/}{Pfreundschuh 等（2022）· Hydronn 巴西近实时卫星降水反演}
\source{R4}{https://arxiv.org/abs/1505.04597}{Ronneberger 等（2015）· U-Net 原论文}
\source{R5}{https://people.eecs.berkeley.edu/~jordan/sail/readings/koenker-bassett.pdf}{Koenker 与 Bassett（1978）· Regression Quantiles}
\source{R6}{https://sites.stat.washington.edu/raftery/Research/PDF/Gneiting2007jasa.pdf}{Gneiting 与 Raftery（2007）· 严格适当评分规则}
\source{R7}{https://ftp.cs.wisc.edu/machine-learning/shavlik-group/davis.icml06.pdf}{Davis 与 Goadrich（2006）· 精确率—召回率与 ROC}
\medskip
\hyperlink{nav}{返回学习路径 $\rightarrow$}

\end{document}
